\addtocontents{toc}{\protect\Needspace{3\baselineskip}}
\originalchapter{2005 年历史测验与解答核对}
B. Zwiebach, 2005 年 3 月 16 日, 限时 60 分钟, 允许个人两页笔记. 原卷七题与五页官方解答均实际取得. 本稿重新推导并与官方答案核对; 尺度商空间原点的拓扑说明为独立补充, 不假称官方证明已处理这一点.

\src{8.251 Historical2005 Q1, 10 分.}
\begin{exercise}
四维时空中两块静止平行 D2 膜沿 $x^1,x^2$, 位于 $x^3=0$ 与 $x^3=a>0$. 弦从第二膜（$\sigma=0$）伸到第一膜（$\sigma=\sigma_1$）. 绘图并逐端点列 $X^0,X^1,X^2,X^3$ 的八个自由或指定值的 Dirichlet 条件.
\end{exercise}
\begin{solution}
两端的 $X^0,X^1,X^2$ 均自由, 即对应三个 $P^\sigma_\mu=0$. 法向分别为 $X^3(t,0)=a$, $X^3(t,\sigma_1)=0$, 变分为零. 用 $P^\sigma$ 表达适用于尚未选共形规范的情形, 共形规范下自由方向等价于 $X'=0$.
\begin{center}
\begin{tikzpicture}[scale=.8]
 \draw[thick](-2,0)--(2,0)node[right]{$x^3=0$};
 \draw[thick](-2,1.6)--(2,1.6)node[right]{$x^3=a$};
 \draw[very thick,structurecolor](.6,1.6)..controls(-.5,1)and(.8,.5)..(-.4,0);
 \node[above]at(.6,1.6){$\sigma=0$};
 \node[below]at(-.4,0){$\sigma_1$};
\end{tikzpicture}
\end{center}
这是端点几何示意, 不声称所画任意曲线是具体动力学解.
\end{solution}

\src{8.251 Historical2005 Q2, (a)--(b), 10 分.}
\begin{exercise}
对复平面分别只施加一个识别:
(a) $z\sim z+\ii$;
(b) $z\sim2z$.
给简单基本域并描述所得空间.
\end{exercise}
\begin{solution}
(a) $z=x+\ii y$, 取 $x\in\mathbb R,0\leq y<1$, 上下边识别, 所得 $\mathbb R\times S^1$ 是无限圆柱.
(b) 在去掉原点的平面, 取 $1\leq|z|<2$, 两边圆按同一辐角配对. 对数半径 $\ln|z|$ 周期为 $\ln2$, 辐角周期 $2\pi$, 所以 $\mathbb C^*/2^{\mathbb Z}$ 是环面.
原官方答案也提示技术上应包含 $z=0$, 本稿进一步说明它不能作为普通平滑环面上的一个正常点: 原点被所有缩放固定, 任意含原点的饱和开邻域都会包含所有非零缩放轨道, 故它无法与其他轨道用不相交开邻域分离. 全平面的商有非 Hausdorff 问题, “所得环面”应明确限定去掉原点的部分. 基本域可另加单独的原点代表, 但这不修复该拓扑.
\end{solution}

\src{8.251 Historical2005 Q3, 15 分.}
\begin{exercise}
三维电磁 Chern--Simons 作用量为
$S=\int\dd t\dd^2x(A_0F_{12}+A_1F_{20}+A_2F_{01})$,
$F_{\mu\nu}=\partial_\mu A_\nu-\partial_\nu A_\mu$.
只对 $A_0$ 变分, 忽略边界项, 用场强写方程.
\end{exercise}
\begin{solution}
$F_{12}$ 不含 $A_0$, $\delta F_{20}=\partial_2\delta A_0$,
$\delta F_{01}=-\partial_1\delta A_0$, 因而
\[
 \delta S=\int\dd^3x\,[\delta A_0F_{12}
 +A_1\partial_2\delta A_0-A_2\partial_1\delta A_0]
 =2\int\dd^3x\,\delta A_0F_{12}.
\]
任意变分给 $F_{12}=0$. 系数 2 来自直接项与经积分分部的另外两项, 不能只变分显式的 $A_0F_{12}$ 而把 $F$ 内的 $A_0$ 漏掉.
\end{solution}

\src{8.251 Historical2005 Q4, (a)--(b), 10 分.}
\begin{exercise}
原自然单位宇宙直弦透镜式为 $\delta=8\pi GT_0$, 亏角以弧度计.
(a) 恢复 $c,\hbar$.
(b) 若 $\delta=0.5$ 角秒, $G=6.7\times10^{-11}\,\mathrm{m^3kg^{-1}s^{-2}}$, $c=3\times10^8\,\mathrm{m\,s^{-1}}$, 求弦线质量密度, 单位 $\mathrm{kg\,m^{-1}}$. 一度为 60 角分, 一角分为 60 角秒.
\end{exercise}
\begin{solution}
$[G]=L^3/(MT^2)$, $[T_0]=ML/T^2$, 故 $GT_0$ 的量纲为 $c^4$, 不需 $\hbar$. 有 $\delta=8\pi GT_0/c^4=8\pi G\mu_0/c^2$.
将角度变为 $\delta=0.5\pi/(180\cdot3600)$, 得
$\mu_0=\delta c^2/(8\pi G)\simeq1.30\times10^{20}\,\mathrm{kg\,m^{-1}}$.
官方文字有一处将 $\mu_0=T_0/c^2$ 排成乘 $c^2$, 其实际计算用的是正确除法, 本稿统一修正.
\end{solution}

\src{8.251 Historical2005 Q5, (a)--(b), 20 分.}
\begin{exercise}
长度 $\ell$、能量 $E$ 的自由开弦以角速度 $\omega$ 刚性旋转, $\omega\ell/2=c$, $\ell=2E/(\pi T_0)$.
(a) 在离中心 $r$ 处长度 $\dd r$ 的小段, 用 $\omega,T_0,r,\dd r,c$ 给相对论动量大小 $\dd p$ 和角动量大小 $\dd J$.
(b) 积分给整弦 $J(E,T_0)$; 可用 $\int_0^1x^2/\sqrt{1-x^2}\dd x=\pi/4$.
\end{exercise}
\begin{solution}
小段静止质量 $T_0\dd r/c^2$, 速度 $\omega r$, 所以
\[
 \dd p=\frac{T_0\omega r\dd r}{c^2\sqrt{1-\omega^2r^2/c^2}},
 \qquad\dd J=r\dd p.
\]
两半角动量同向, 积分后
\[
 J=\frac{2T_0\omega}{c^2}\int_0^{c/\omega}
 \frac{r^2\dd r}{\sqrt{1-\omega^2r^2/c^2}}
 =\frac{2T_0c}{\omega^2}\frac\pi4.
\]
由 $\omega=\pi T_0c/E$ 得 $J=E^2/(2\pi T_0c)$.
端点密度发散可积; 不能因为密度在端点无穷就断言整弦总量无穷.
\end{solution}

\src{8.251 Historical2005 Q6, (a)--(c), 25 分.}
\begin{exercise}
闭弦在均匀能量参数下
$\boldsymbol X=\frac12[\boldsymbol F(u)+\boldsymbol G(v)]$,
$u=ct+\sigma,v=ct-\sigma$, $\sigma\sim\sigma+L$, $L=E/T_0$.
(a) 由周期性联系 $\boldsymbol F,\boldsymbol G$ 的非周期部分.
(b) 写 $\boldsymbol F=\boldsymbol f+\boldsymbol\alpha u$,
$\boldsymbol G=\boldsymbol g+\boldsymbol\beta v$, $\boldsymbol f,\boldsymbol g$ 严格 $L$ 周期, 求 $\alpha,\beta$ 关系并化简 $\boldsymbol X$.
(c) 动量密度为 $\boldsymbol P^t=(T_0/c^2)\partial_t\boldsymbol X$, 求整弦动量.
\end{exercise}
\begin{solution}
$\boldsymbol X(t,\sigma+L)=\boldsymbol X(t,\sigma)$ 给
$\boldsymbol F(u+L)-\boldsymbol F(u)=\boldsymbol G(v)-\boldsymbol G(v-L)$.
代入分解得 $\boldsymbol\alpha L=\boldsymbol\beta L$, 因而
\[
 \boldsymbol X=\boldsymbol\alpha ct+
 \tfrac12[\boldsymbol f(ct+\sigma)+\boldsymbol g(ct-\sigma)].
\]
对时间求导并积分, 两个周期导数各自整圈积分为零, 故
$\boldsymbol p=(T_0L/c)\boldsymbol\alpha=(E/c)\boldsymbol\alpha$.
这是质心速度 $\boldsymbol v_{\rm cm}=c\boldsymbol\alpha$ 与
$\boldsymbol p=E\boldsymbol v_{\rm cm}/c^2$ 的一致性检验.
\end{solution}

\src{8.251 Historical2005 Q7, 10 分.}
\begin{exercise}
闭弦沿一个椭圆静止拉开, 再以零初速度释放. 圆形闭弦会在后来收缩至零大小, 椭圆是否也会? 给精确理由, 不必做繁复积分.
\end{exercise}
\begin{solution}
会. 初始静止使能量参数恰为弧长, 设椭圆周长 $L$, 中心为原点, 初始弧长曲线为 $\boldsymbol f(\sigma)$. 中心反演对称给
$\boldsymbol f(\sigma+L/2)=-\boldsymbol f(\sigma)$.
零初速度解是
$\boldsymbol X(t,\sigma)=\frac12[\boldsymbol f(\sigma+ct)+\boldsymbol f(\sigma-ct)]$.
在 $t=L/(4c)$, 两个参数相差 $L/2$, 因此对每个 $\sigma$ 都相消, 整弦到达原点. 此理由也适用于光滑、弧长参数一致且中心反演对称的其他初始闭合曲线; 不能用“椭圆各处曲率不同”来否定这个全局解性质.
\end{solution}

